\subsection{PRIMITIVE ELEMENTS OF A BIGEBRA}

Let E be a bigebra (Algebra, Chapter III, § 11, no. 4), c its coproduct, ε its counit and 1 its unit element. As ε(1) = 1 and c(1) = 1 ⊗ 1, the results of the previous no. can be applied with u = 1. The 1-primitive elements of E (no. 1, Definition 1) are simply called primitive (cf.~Algebra, Chapter III; § 11, no. 8), that is the elements x of E such that

\[
c (x) = x \otimes 1 + 1 \otimes x.\tag{5}
\]

We write simply \(\pi\) , \(\eta\) , P(E), \(c^{+}\) instead of \(\pi_{1}\) , \(\eta_{1}\) , P \(_{1}\) (E), \(c_{1}^{+}\) .

PROPOSITION 4. The set \(\mathrm{P(E)}\) of primitive elements of \(\mathbf{E}\) is a Lie subalgebra of \(\mathbf{E}\) .

If \(x, y\) are in \(\mathrm{P(E)}\) , then

\[
\begin{array}{r l} c (x y) & = c (x) c (y) = (x \otimes 1 + 1 \otimes x) (y \otimes 1 + 1 \otimes y) \\ & = x y \otimes 1 + 1 \otimes x y + x \otimes y + y \otimes x, \end{array}
\]

whence

\[
c ([ x, y ]) = [ x, y ] \otimes 1 + 1 \otimes [ x, y ].
\]

PROPOSITION 5. Let \(f: \mathrm{E} \to \mathrm{E}'\) be a bigebra morphism. If \(x\) is a primitive element of \(\mathrm{E}\) , then \(f(x)\) is a primitive element of \(\mathrm{E}'\) and the restriction of \(f\) to \(\mathrm{P(E)}\) is a Lie algebra homomorphism \(\mathrm{P}(f): \mathrm{P(E)} \to \mathrm{P(E')}\) .

Let \(c\) (resp. \(c'\) ) be the coproduct of E (resp. E'). Since \(f\) is a cogebra morphism \(c' \circ f = (f \otimes f) \circ c\) , whence

\[
\begin{array}{r l} c ^ {\prime} (f (x)) & = (f \otimes f) (c (x)) = (f \otimes f) (x \otimes 1 + 1 \otimes x) \\ & = f (x) \otimes 1 + 1 \otimes f (x), \end{array}
\]

for \(x\) primitive. Hence \(f\) maps \(\mathrm{P(E)}\) into \(\mathrm{P(E')}\) and \(f([x,y]) = [f(x),f(y)]\) since \(f\) is an algebra homomorphism.

Remarks. (1) Let \(p\) be a prime number such that \(p\) . \(1 = 0\) in \(\mathbf{K}\) . The binomial formula and the congruences \(\binom{p}{i} \equiv 0 \pmod{p}\) for \(1 \leqslant i \leqslant p - 1\) imply that

P(E) is stable under the mapping \(x \mapsto x^p\) .

\begin{enumerate}
\def\labelenumi{(\arabic{enumi})}
\setcounter{enumi}{1}
\tightlist
\item
  By definition, the diagram
\end{enumerate}

\[
0 \longrightarrow \mathrm{P(E)} \longrightarrow \mathrm {E^ {+}} \stackrel {c ^ {+}} {\longrightarrow} \mathrm {E^ {+}} \otimes \mathrm {E^ {+}}
\]

is an exact sequence. If \(\mathbf{K}'\) is a commutative ring and \(\rho: \mathbf{K} \to \mathbf{K}'\) a ring homomorphism, \(\rho^*(\mathrm{E}) = \mathrm{E} \otimes_{\mathbb{K}} \mathbf{K}'\) is a \(\mathbf{K}'\) -bigebra and the inclusion \(\mathrm{P}(\mathrm{E}) \to \mathrm{E}\) defines a homomorphism of Lie \(\mathbf{K}'\) -algebras

\[
\alpha \colon \mathrm{P} (\mathrm{E}) \otimes_ {\mathrm{K}} \mathrm{K} ^ {\prime} \rightarrow \mathrm{P} (\mathrm{E} \otimes_ {\mathrm{K}} \mathrm{K} ^ {\prime}).
\]

If \(\mathbf{K}'\) is flat over \(\mathbf{K}\) (Commutative Algebra, Chapter I, §2, no. 3, Definition 2), it follows from loc. cit. that the diagram

\[
0 \longrightarrow \mathrm{P} (\mathrm{E}) \otimes_ {\mathrm{K}} \mathrm{K} ^ {\prime} \longrightarrow \mathrm{E} ^ {+} \otimes_ {\mathrm{K}} \mathrm{K} ^ {\prime} \xrightarrow {c ^ {+} \otimes_ {\mathrm{K}} \mathrm{Id} _ {\mathrm{K} ^ {\prime}}} \left(\mathrm{E} ^ {+} \otimes_ {\mathrm{K}} \mathrm{K} ^ {\prime}\right) \otimes_ {\mathrm{K} ^ {\prime}} \left(\mathrm{E} ^ {+} \otimes_ {\mathrm{K}} \mathrm{K} ^ {\prime}\right)
\]

is an exact sequence, which implies that \(\alpha\) is an isomorphism.

